\documentclass[12pt,a4paper,landscape]{article}
\usepackage[utf8]{inputenc}
\usepackage[ukrainian]{babel}
\usepackage{graphicx}
\usepackage{multirow}
\usepackage{array}
\usepackage{booktabs}
\usepackage{longtable}
\usepackage{xcolor}
\usepackage{colortbl}
\usepackage{geometry}
\geometry{left=2cm,right=2cm,top=2cm,bottom=2cm}

\begin{document}

\section*{Додаток А. Детальний аналіз тематичних кластерів досліджень з навчання інформаційній безпеці}

\begin{longtable}{|p{1cm}|p{3.5cm}|p{5cm}|p{4cm}|p{4cm}|p{5cm}|}
\caption{Характеристика тематичних кластерів за результатами VOSviewer аналізу}\\
\hline
\rowcolor{gray!20}
\textbf{№} & \textbf{Назва кластера} & \textbf{Ключові терміни} & \textbf{Характеристики} & \textbf{Новітні тренди} & \textbf{Наукові імплікації} \\
\hline
\endfirsthead

\multicolumn{6}{c}%
{\tablename\ \thetable\ -- \textit{Продовження}} \\
\hline
\rowcolor{gray!20}
\textbf{№} & \textbf{Назва кластера} & \textbf{Ключові терміни} & \textbf{Характеристики} & \textbf{Новітні тренди} & \textbf{Наукові імплікації} \\
\hline
\endhead

\hline
\endfoot

\hline
\endlastfoot

\cellcolor{red!20}1 & 
\textbf{Інформаційні технології та їхній вплив на суспільство} & 
\begin{itemize}
    \item artificial intelligence (85 зв'язків, 18 появ)
    \item machine learning (середній рік 2022.1)
    \item e-learning (233 зв'язки, 52 появи)
    \item data privacy
    \item virtual reality
    \item social media
    \item covid-19 (середній рік 2022.5)
\end{itemize} & 
\begin{itemize}
    \item Найбільш загальний кластер
    \item Охоплює сучасні технології
    \item Високий середній рік публікацій (2021-2022)
    \item Висока нормалізована цитованість
\end{itemize} & 
\begin{itemize}
    \item Інтеграція ШІ в онлайн-освіту
    \item AI-асистенти для навчання
    \item Адаптивне персоналізоване навчання
    \item VR/AR технології в освіті
    \item Відповідь на виклики пандемії
\end{itemize} & 
\begin{itemize}
    \item Розробка цифрових освітніх систем
    \item Оцінка ефективності AI-платформ
    \item Дослідження впливу технологій на навчання
    \item Етичні питання використання ШІ в освіті
\end{itemize} \\
\hline

\cellcolor{green!20}2 & 
\textbf{Освіта та обізнаність у сфері інформаційної безпеки} & 
\begin{itemize}
    \item education (356 зв'язків, 72 появи)
    \item information security education (45 зв'язків, 12 появ)
    \item security of data (425 зв'язків, 99 появ)
    \item authentication
    \item security awareness
    \item information systems
\end{itemize} & 
\begin{itemize}
    \item Фокус на формальній освіті
    \item Високі показники зв'язків
    \item Фундаментальні теми
    \item Середній рік публікацій 2017-2018
\end{itemize} & 
\begin{itemize}
    \item Розвиток цифрової грамотності
    \item Стандарти кібергігієни
    \item Інтеграція безпеки в навчальні програми
    \item Формування культури безпеки
\end{itemize} & 
\begin{itemize}
    \item Підготовка викладачів
    \item Розробка курсів з ІБ
    \item Політика безпечного навчання
    \item Методики оцінки обізнаності
\end{itemize} \\
\hline

\cellcolor{blue!20}3 & 
\textbf{Навчання кібербезпеці в освітніх закладах} & 
\begin{itemize}
    \item cyber security (1008 зв'язків, 195 появ)
    \item cybersecurity education (209 зв'язків, 49 появ)
    \item capture the flag (CTF)
    \item game-based learning
    \item adversarial machine learning (середній рік 2024.3)
    \item k-12 education
\end{itemize} & 
\begin{itemize}
    \item Найпотужніший кластер за зв'язками
    \item Активні сучасні дослідження
    \item Інноваційні методи навчання
    \item Орієнтація на школярів
\end{itemize} & 
\begin{itemize}
    \item CTF-змагання як метод навчання
    \item Симуляції кібератак
    \item Практична hands-on підготовка
    \item Включення ML/AI в безпеку
    \item Геймифіковане навчання
\end{itemize} & 
\begin{itemize}
    \item Методики практичного навчання
    \item Розробка навчальних симуляторів
    \item Адаптація для різних вікових груп
    \item Оцінка ефективності ігрових методів
\end{itemize} \\
\hline

\cellcolor{yellow!20}4 & 
\textbf{Тренування персоналу в галузі кібербезпеки} & 
\begin{itemize}
    \item cybersecurity (862 зв'язки, 195 появ)
    \item computer crime (182 зв'язки, 29 появ)
    \item network security (281 зв'язок, 52 появи)
    \item personnel training
    \item serious games
    \item gamification
\end{itemize} & 
\begin{itemize}
    \item Фокус на професійній підготовці
    \item Висока цитованість (14.3 для computer crime)
    \item Практична спрямованість
    \item Використання серйозних ігор
\end{itemize} & 
\begin{itemize}
    \item Інцидент-менеджмент у школах
    \item Викладання через симулятори
    \item Серйозні ігри для тренування
    \item Сценарії реагування на загрози
\end{itemize} & 
\begin{itemize}
    \item Інтеграція кіберзахисту в лабораторії
    \item STEM-курси з кібербезпеки
    \item Підвищення кваліфікації вчителів
    \item Розробка тренінгових програм
\end{itemize} \\
\hline

\cellcolor{violet!20}5 & 
\textbf{Освіта в комп'ютерних науках та інженерії} & 
\begin{itemize}
    \item students (1008 зв'язків, 190 появ)
    \item computer science education
    \item curricula (433 зв'язки, 76 появ)
    \item engineering education (402 зв'язки, 73 появи)
    \item teaching
    \item k-12, middle school
\end{itemize} & 
\begin{itemize}
    \item Студентоцентрований підхід
    \item Академічна спрямованість
    \item Фокус на навчальних програмах
    \item Охоплення різних освітніх рівнів
\end{itemize} & 
\begin{itemize}
    \item Розвиток інформатики в школах
    \item STEAM-підходи
    \item Інтеграція з математикою та науками
    \item Проектне навчання
\end{itemize} & 
\begin{itemize}
    \item Впровадження інформатики в програми
    \item Оцінка ефективності навчання
    \item Розробка стандартів освіти
    \item Підготовка навчальних матеріалів
\end{itemize} \\
\hline

\end{longtable}

\vspace{1cm}

\section*{Основні висновки з кластерного аналізу:}

\begin{enumerate}
    \item \textbf{Домінуючі теми:} Кібербезпека (``cyber security'', ``cybersecurity'') є центральною темою дослідження з найвищими показниками зв'язків (1008) та кількості появ (195) у кількох кластерах.
    
    \item \textbf{Взаємозв'язок тематик:} Існує сильний зв'язок між темами кібербезпеки та освітою, що видно з багатьох термінів-мостів між кластерами (``cyber-security educations'', ``cybersecurity education'', ``information security education'').
    
    \item \textbf{Часова актуальність:} Дослідження в галузі кібербезпеки, штучного інтелекту та машинного навчання є найбільш актуальними (середні роки публікації 2021-2024), що свідчить про активний розвиток цих напрямів.
    
    \item \textbf{Інновації в навчанні:} Використання ``gamification'', ``serious games'', ``capture the flag'' та ``game-based learning'' вказує на пошук нових, інтерактивних підходів до навчання кібербезпеці.
    
    \item \textbf{Фокус на аудиторії:} Значна увага приділяється різним групам учнів: ``students'' (190 появ), ``high school students'', ``k-12 education'', ``middle school'', ``primary schools'', ``secondary schools'', ``teachers'', ``personnel'', що свідчить про диференційований підхід до навчання.
    
    \item \textbf{Міждисциплінарність:} Перетин технічних (cybersecurity, network security), педагогічних (education, teaching) та соціальних (awareness, privacy) аспектів демонструє комплексний характер проблеми.
\end{enumerate}

\end{document}